\ficha{Abreu (2014), A disputa monetária na Primeira República}{abreu2014}

\begin{identificacao}
ABREU, Cristiano Addario de. \emph{A disputa monetária na Primeira República
(1890-1906). Entre papelistas e metalistas: a moeda como projeção e resultado do
real}. 2014. Dissertação (Mestrado em História Econômica), Faculdade de
Filosofia, Letras e Ciências Humanas, Universidade de São Paulo, São Paulo, 2014.
Orientador: Prof. Dr. Alexandre de Freitas Barbosa. DOI:
10.11606/D.8.2014.tde-31032015-113729.

Dissertação de mestrado, 117 páginas, em português. Lida integralmente, na versão
corrigida depositada.
\end{identificacao}

\bloco{Problema e tese}
A pergunta é qual moeda foi desejada e qual foi possível no Brasil entre a
proclamação da República e o Convênio de Taubaté, e que grupos sustentaram cada
resposta. A tese é que a política monetária republicana oscila como pêndulo entre
os dois polos herdados do Império, o papelista de Rui Barbosa em 1890 e o
metalista de Campos Sales e Joaquim Murtinho em 1898, e que 1906 não é vitória de
um polo, mas síntese pragmática que Abreu batiza de ``metalismo possível'': a
Caixa de Conversão permitia ampliar o meio circulante com lastro em cambiais e ao
mesmo tempo manter taxa estável e relativamente valorizada (p. 59).

\bloco{Argumento}
A dissertação inteira propõe ler a moeda como sintoma do deslocamento do eixo
dinâmico da economia brasileira, do setor exportador para o mercado interno. Cada
arranjo monetário testado entre 1889 e 1906 é tratado como expressão de quanto o
Brasil conseguia, em cada conjuntura de balanço de pagamentos, aproximar-se do
padrão-ouro sem asfixiar a monetização interna. O quadro teórico vem de Arrighi,
Braudel, Rangel e Furtado, não da teoria monetária.

\begin{enumerate}[leftmargin=*,nosep]
\item As categorias são importadas. O confronto brasileiro é ``reprise com
variações'' da batalha inglesa do Bullion Report de 1810, em que Ricardo estava do
lado metalista ou bulionista, e que prossegue até o Bank Charter Act de 1844
(p. 29, citando Gustavo Franco). Abreu lembra que em quase metade do século da
\emph{pax britannica} o padrão-ouro esteve ausente do próprio centro, e que entre
1846 e 1906 o câmbio brasileiro esteve perto da paridade em apenas 60 de 720
meses, 8,3\% do tempo (p. 30).

\item A taxonomia é curta e operacional. Metalistas defendem conversibilidade em
ouro com emissão controlada pelo Estado; papelistas defendem pluralidade de
bancos emissores e papel-moeda inconversível lastreado em títulos da dívida
pública; o capital financeiro externo fica com os primeiros, porque a
conversibilidade assegura o repatriamento (p. 29). Metalistas se preocupam com a
imagem e a estabilidade da moeda, papelistas com a escassez de crédito que
travaria o desenvolvimento (pp. 30-31). O indicador da política monetária, para
os papelistas, é a taxa de juros, não o câmbio (p. 39).

\chave{por mais que muitos considerassem os metalistas dogmáticos, tinham eles
então mais estofo teórico, estando os papelistas calcados mais intuitivamente do
que doutrinariamente}{p. 38}

\item O século XIX é pendular, com pontos marcados: Alves Branco (1844), Itaboraí
em 1853, a pluralidade emissora de Souza Franco em 1857, a crise do Souto em 1864
que obriga os próprios saquaremas metalistas ao curso forçado, e a guerra do
Paraguai (pp. 31-32).

\item A contribuição própria do autor é situar o Encilhamento na contradição entre
a reserva bancária fracionária e o padrão-ouro (pp. 37-38). A ``cláusula de
exceção'' que autorizava o curso forçado em caso de corrida era, na prática,
cláusula de regra (p. 37). A alavancagem começa no triplo com conversibilidade no
gabinete Ouro Preto, passa ao dobro sem conversibilidade com Rui Barbosa e volta
ao triplo sem conversibilidade na fusão que cria o BREUB (p. 43). O vaivém emissor
do próprio ministério, somado à crise Baring e à queda do preço do café, produziu
a incerteza que alimentou a crise.

\item O funding loan de 1898 é reação metalista e subordinação da política interna
ao serviço da dívida. Negociado com os Rothschild, com hipoteca das rendas da
Alfândega do Rio, obrigava a retirar de circulação papel-moeda equivalente ao
empréstimo ao câmbio de 18 pence (p. 46). Campos Sales foi pessoalmente a Londres
antes de assumir e cumpriu o acordo de forma literal; Murtinho incinerou papel e
restabeleceu a quota-ouro nas tarifas alfandegárias (p. 47). O resultado foi
recessivo até 1905.

\item A estabilização de Murtinho fica longe da meta metalista. O câmbio ficou em
torno de 14 pence por mil-réis sob o \emph{funding scheme}, contra o par legal de
27 fixado em 1846, e o ponto mais baixo de 1891 tinha sido 11 pence (p. 75). A
apreciação é atribuída mais à melhora do balanço de pagamentos, com a moratória e
o \emph{boom} da borracha, do que à contração monetária.

\chave{não sabemos até onde a apreciação cambial se deve à contração monetária ou
a fatores associados ao balanço de pagamentos, como exportações de borracha e a
retomada das entradas de capital}{p. 75}

\item Os grupos não se alinham pela geografia. Três presidentes paulistas
seguidos governaram contra a lavoura cafeeira porque se apoiavam no grande
capital cafeeiro, diversificado em ferrovias, bancos e intermediação, que preferia
câmbio alto e estável; o apelo baixista vinha do pequeno e médio capital
(pp. 51-53, com Saes, Kugelmas e Perissinotto). Os cafeicultores teriam sido, na
primeira década, mais criaturas do que criadores da política papelista (p. 49).

\item 1906 é ponto de mutação. Diante da safra-monstro, os estados produtores
intervêm, e a cláusula mais controversa do Convênio é a que pedia ao governo
federal uma caixa de conversão que fixasse a taxa do mil-réis (p. 58). Abreu adota
a avaliação de Holloway de que o convênio nunca teve força de lei e não foi
executado como planejado, e diz que o termo valorização é errôneo, pois o que se
conseguiu foi evitar uma desvalorização (p. 58). A Caixa, inovação vinda da
Argentina, funciona como antepassado de banco central e sustentou crescimento não
inflacionário até 1914 (p. 77).
\end{enumerate}

\bloco{Conceitos e definições operacionais}
Metalismo e papelismo, como acima (p. 29). \emph{Real bills doctrine}: base
doutrinária dos papelistas, a noção de que emissão em desconto de duplicata
legítima jamais poderia ser inflacionária (pp. 38-39). Reserva bancária
fracionária: o que a época chamava de encaixe bancário ou alavancagem (p. 37).
``Metalismo possível'': categoria do próprio autor para 1906, emissão contra
cambiais com taxa administrada, que concilia expansão do numerário e compromisso
cambial (pp. 59, 77). Baixa do câmbio: o autor avisa duas vezes que o vocabulário
de época inverte o atual, e que baixa do câmbio é desvalorização do mil-réis
(pp. 35-36, 80).

\bloco{Evidência e método}
Não há série nova, teste nem contrafactual. O capítulo 3 usa fontes primárias
diretas: o discurso de Rui Barbosa no Senado de 3 de novembro de 1891, publicado
como \emph{O Papel e a Baixa do Câmbio}, e três relatórios ministeriais de
Murtinho, o da Indústria, Viação e Obras Públicas de 1897 e os da Fazenda de 1899
e 1901, pela coletânea \emph{Ideias Econômicas de Joaquim Murtinho} (p. 79).
Aparecem ainda \emph{A Presidência Campos Salles} de Alcindo Guanabara (1902), um
parecer de Murtinho no Senado em 1896 e o manifesto de Manuel Vitorino de 1898. Os
capítulos 1 e 2 são historiográficos, com Furtado, Peláez e Suzigan, Gustavo
Franco, Holloway, Delfim Netto, Levy, Saes, Kugelmas, Perissinotto e Palazzo. As
duas tabelas relevantes são reproduzidas de terceiros: preço do café importado
nos Estados Unidos, de Delfim Netto (p. 13), e balanço de pagamentos 1889-1900, de
Gustavo Franco (pp. 74-75).

\bloco{Agentes e vertentes}
Papelistas: Rui Barbosa, cuja reforma de 17 de janeiro de 1890 é lida como
importação declarada do \emph{National Banking Act} e cuja moeda monetizada com
juros regressivos até zero é aproximada do \emph{greenback} de Lincoln (pp. 39-40,
87). Rui não recusa o padrão-ouro como meta, apenas inverte a ordem causal: é a
prosperidade que permite a conversibilidade, e não o contrário (p. 87).
Metalistas: Murtinho, cujo liberalismo se apoia em darwinismo social e em
argumento racial explícito contra tomar os Estados Unidos como modelo (p. 104);
aceita em tese o papel inconversível que financie ``trabalho produtivo'', mas nega
ao Brasil essa capacidade (pp. 92-93, 97). Campos Sales é o executor literal do
funding. Alcindo Guanabara entra só como fonte de época sobre o contrato do Banco
Nacional (p. 36). Vieira Souto é o caso mais útil aqui: professor da Escola
Politécnica, publicou no \emph{Correio da Manhã}, entre novembro e dezembro de
1901, a série depois reunida em \emph{O Último Relatório da Fazenda}, contra a
equação cambial de Murtinho, sustentando que o câmbio é regido pelos créditos e
débitos recíprocos e que o papel-moeda é causa apenas indireta (p. 76). Do lado da
teoria monetária do câmbio, Abreu lista Antônio Carlos Ribeiro de Andrada, Carlos
Inglês de Sousa, Pandiá Calógeras e Ramalho Ortigão (p. 75). Em Taubaté: Jorge
Tibiriçá, Nilo Peçanha, Francisco Sales, Cândido Rodrigues, Albuquerque Lins e
Augusto Ramos, com antecedentes em Quintino Bocaiuva (1902) e Alexandre Siciliano
(1903). Os Rothschild são contraparte do funding, sem estudo próprio.

\bloco{Críticas e limites}
O texto não menciona Leopoldo de Bulhões, David Campista nem Serzedello Corrêa:
quem procurar aqui a genealogia dos agentes de 1906 encontra Rui Barbosa e
Murtinho, e mais nada. A pesquisa de imprensa é indireta, e os únicos jornais
citados são o \emph{Correio da Manhã}, pela série de Vieira Souto, e uma
referência genérica à imprensa fluminense que atacava Rui em 1891 (p. 88). A tese
sobre 1906 apoia-se quase toda em Holloway, Delfim Netto e Saes. A Caixa de
Conversão não é estudada em seu funcionamento: não há limite de emissão, fundo de
garantia, nível da taxa fixada, agência de Londres nem suspensão, e 1914 comparece
só como data em que a guerra encerra o arranjo (p. 111). A comparação com a Caja
argentina é afirmada em uma linha, sem apoio (pp. 60, 77). Murtinho é lido apenas
pelos relatórios, sem confronto com a execução orçamentária. A moldura arrighiana
faz trabalho retórico, e o vínculo entre ciclo sistêmico e escolha monetária
brasileira é asseverado, não demonstrado.

\begin{dialogo}
\item Sobrevida do vocabulário. Procurar, em 1906, ``metalista'' e ``papelista''
aplicados a contemporâneos, e não à memória do Encilhamento. Abreu sustenta que
1906 é síntese dos dois polos. Se o vocabulário de época já não classifica os
debatedores de 1906, isso confirma a decisão de usar ortodoxo e expansionista como
eixo analítico.

\item Teoria monetária contra balanço de pagamentos. Procurar, nos editoriais
sobre a taxa e sobre o limite de emissão, se a variação cambial é atribuída ao
volume de emissão ou ao saldo externo, safra e entrada de capitais. Abreu diz que
a primeira explicação dominava e nomeia Vieira Souto como exceção publicada no
\emph{Correio da Manhã} em 1901 (pp. 75-76). Testa se em 1906 a exceção virou
posição corrente, e se aquele jornal manteve a linha.

\item \emph{Real bills} e necessidades do comércio. Procurar o argumento de que a
emissão contra cambiais não é inflacionária porque atende às necessidades do
comércio, e a menção à taxa de juros como indicador, núcleo doutrinário papelista
segundo Abreu (pp. 38-39). Testa se a defesa da Caixa reaproveita a doutrina
papelista ou argumenta apenas pelo lastro.

\item Quem fala pela lavoura. Procurar se os jornais tratam a lavoura como bloco
ou distinguem grande capital cafeeiro, com interesses em ferrovias e bancos, do
pequeno e médio produtor. Abreu faz dessa clivagem a explicação para governos
paulistas contrariarem cafeicultores (pp. 51-53). Testa se ela aparece no discurso
público ou é só construção historiográfica.

\item Herança do funding. Procurar menções ao funding loan, aos Rothschild, à
alfândega hipotecada e à incineração de papel de Murtinho como argumento a favor
ou contra a Caixa. Abreu trata 1898 como o trauma que organiza a virada de 1906.
Testa se os jornais mobilizam 1898 como advertência ou como modelo de
credibilidade.
\end{dialogo}

\usonocapitulo{Parte III, como espinha dorsal cronológica. Serve para três
afirmações. Primeira, que metalismo e papelismo chegam ao Brasil como importação
das categorias inglesas do debate bulionista e do Bank Charter Act, com taxonomia
de época datável (p. 29). Segunda, que o pêndulo entre os dois polos percorre o
século XIX e a primeira década republicana, e que em 1906 a clivagem já não separa
os campos, o que sustenta o uso do eixo ortodoxo contra expansionista no período
da Caixa e a reserva de metalismo e papelismo ao vocabulário de época. Terceira,
que a Caixa de Conversão foi lida por parte da elite como conciliação entre
expansão do numerário e compromisso cambial, o ``metalismo possível'' de Abreu
(pp. 59, 77), hipótese que a leitura dos jornais pode confirmar ou negar. Alimenta
também a Parte II, com o funding loan de 1898 visto do lado brasileiro, ao lado de
Weller e de Flandreau e Flores. Não serve à Parte IV: a comparação com a Caja
argentina é mencionada sem desenvolvimento.}
